% Pursuing Stacks (cotes 134-2 to 134-8): archive cote 134-7, site page 71 (PDF page 72), typescript page 513.
% typescript. Draft compared: ps6.tex:3526-3605.
% AI-transcribed by Claude Opus 5.5 (claude-opus-5-5, adaptive thinking, effort high), Message Batches API, batch ps4 (msgbatch_012Aoc1NuskATPsyBeKSYK7J), 2026-11-12; not yet checked by eye.
\scanpage{134-7/71}
\mspage{513}

Truth to tell, these generalities are more interesting still \xout{in} for
an argument X in a category like $\hand{\Delta}^{\wedge}{}_{/Y} \simeq (\hand{\Delta}_{/Y})^{\wedge}$ , where Y is an arbitrary
object in $\hand{\Delta}^{\wedge}$, rather than in $\hand{\Delta}^{\wedge}$ itself. In the latter case, X always
admits a pointed structure i.e. a section over the final object (provided
X is non-empty), hence there always exists a splitting for (1), and
hence \inserted{one} for (6); whereas for an object X over Y, there does not always
exist a section over Y, and accordingly it may well happen that the exten%
% NOTE: the right edge of the scan cuts the end of this line and of some lines below
-sion similar to (1) (but taken ``over Y'') does not split, i.e. that the
torsor $W_{/Y}(X)(1)$ is not trivial. The class of this torsor is an element
\xout{in}
\[ (7) \qquad c(X/Y)\hand{\in} H^1(Y, \hand{\widetilde{W}}_{/Y}(X)) \quad , \]
a very interesting invariant indeed, giving rise to the Postnikov inva-%
riants, when Y is one of the $X_n$'s occuring in the Postnikov dévissage
of X . It was while trying to understand the precise relationship between
a cohomology group as in (\unsure{7}), with ``continu\xout{\unsure{o}}ous'' coefficients (by which
I mean to suggest that $\hand{\widetilde{W}}_{/Y}(X)$ is viewed intuitively as a fiberspace over
the ``space'' X, whose fibers are topological \inserted{abelian} groups which are by no means
discrete, but \xout{\illegible{5}}\inserted{got a bunch of non vanishing $\pi_i$'s !),} and a cohomology group
with ``discrete'' coefficients, more accurately with coefficients in a
complex of chains in the category of abelian sheaves over X i.e. over $\hand{\Delta}_{/X}$,
whose homology sheaves are definitely of a ``diserete'' nature\xout{,} (for instanc%
they are locally constant if X$\to$Y is a Kan fibration, and their fibers
are $\underline{Z}$-modules of finite type provided we make a mild finite-type assumpti%
on on the fibers of X$\to$Y), that I got involved in a more general reflec-%
tion on a Kan\hand{-}Dold-Puppe type of relationship, \xout{which} I hope to come back
to \inserted{this} when getting back to the general review o\over{n}{f} linearization begun in part
V of these notes, and abruptly interrupted after section 109, when
getting \xout{involved} caught unwittingly by the enticing mystery of schematiza-%
tion !

Another afterthought to the previous notes, is that I am going to
denote by K , $\hand{\widetilde{K}}$ , the functors
\[ (7) \qquad K\colon \mathrm{Hotab} \longrightarrow \mathrm{Hot} \quad , \quad \hand{\widetilde{K}} : \mathrm{Hotab} \to \mathrm{Hot}^{\hand{\bullet}} \hand{,} \quad \hand{K=\beta\circ\widetilde{K} ,} \]
denoted by $K_{\hand{\pi}}$ , $\hand{\widetilde{K}}_{\hand{\pi}}$ in section 121. This gives a formula such as
\[ (8) \qquad K(\hand{\pi}[n]) = K(\hand{\pi},n) \quad , \]
where $\hand{\pi}[n]$ denotes the chain complex in (Ab) which is $\hand{\pi}$ placed in degree
n . I was intending to complement the \xout{\unsure{latter}} former notation $K_{\hand{\pi}}$ by a
similar notation $K_H$ (where H stands for ``homology''\hand{,} as $\hand{\pi}$ was standing

\margin{71}
